\begin{exercise}[VI.29.E1: Chain length]\label{exer:vi29-01}
What is the length of $\mathfrak p_0\subsetneq\mathfrak p_1\subsetneq\mathfrak p_2$?
\end{exercise}

\begin{exercise}[VI.29.E2: Field]\label{exer:vi29-02}
Compute $\dim k$ for a field $k$.
\end{exercise}

\begin{exercise}[VI.29.E3: Integers]\label{exer:vi29-03}
Compute $\dim\mathbb Z$.
\end{exercise}

\begin{exercise}[VI.29.E4: Affine line]\label{exer:vi29-04}
Compute $\dim k[x]$.
\end{exercise}

\begin{exercise}[VI.29.E5: Affine plane lower bound]\label{exer:vi29-05}
Give a prime chain of length $2$ in $k[x,y]$.
\end{exercise}

\begin{exercise}[VI.29.E6: Height]\label{exer:vi29-06}
Define $\operatorname{ht}(\mathfrak p)$.
\end{exercise}

\begin{exercise}[VI.29.E7: Localization]\label{exer:vi29-07}
Why can localization not increase Krull dimension?
\end{exercise}

\begin{exercise}[VI.29.E8: Quotient]\label{exer:vi29-08}
Why can quotienting not increase dimension?
\end{exercise}

\begin{exercise}[VI.29.E9: Reduction]\label{exer:vi29-09}
Show nilpotents do not change dimension.
\end{exercise}

\begin{exercise}[VI.29.E10: Product]\label{exer:vi29-10}
Compute $\dim(A\times B)$ from $\dim A$ and $\dim B$.
\end{exercise}

\begin{exercise}[VI.29.E11: Dual numbers]\label{exer:vi29-11}
Compute $\dim k[\varepsilon]/(\varepsilon^2)$.
\end{exercise}

\begin{exercise}[VI.29.E12: Reducible curve]\label{exer:vi29-12}
Compute $\dim k[x,y]/(xy)$.
\end{exercise}

\begin{exercise}[VI.29.E13: Height in the plane]\label{exer:vi29-13}
Compute $\operatorname{ht}(x)$ in $k[x,y]$.
\end{exercise}

\begin{exercise}[VI.29.E14: Maximal height]\label{exer:vi29-14}
Compute $\operatorname{ht}(x,y)$ in $k[x,y]$.
\end{exercise}

\begin{exercise}[VI.29.E15: Infinite variables]\label{exer:vi29-15}
Why does $k[x_1,x_2,\ldots]$ have infinite dimension?
\end{exercise}

\begin{exercise}[VI.29.E16: Normalization]\label{exer:vi29-16}
Why does normalization preserve Krull dimension?
\end{exercise}

\begin{exercise}[VI.29.E17: Principal open]\label{exer:vi29-17}
Interpret $\dim A_f$ in terms of $\Spec A$.
\end{exercise}

\begin{exercise}[VI.29.E18: Closed subscheme]\label{exer:vi29-18}
Interpret $\dim(A/I)$ in terms of $V(I)$.
\end{exercise}

\begin{exercise}[VI.29.E19: Specialization]\label{exer:vi29-19}
When is $\mathfrak q$ a specialization of $\mathfrak p$ in $\Spec A$?
\end{exercise}

\begin{exercise}[VI.29.E20: Irreducible chain]\label{exer:vi29-20}
Translate $(0)\subsetneq(x)\subsetneq(x,y)$ into closed subsets of $\Spec k[x,y]$.
\end{exercise}

\begin{exercise}[VI.29.E21: Polynomial dimension]\label{exer:vi29-21}
State $\dim k[x_1,\ldots,x_n]$.
\end{exercise}

\begin{exercise}[VI.29.E22: Artinian ring]\label{exer:vi29-22}
What is the dimension of a nonzero Artinian ring?
\end{exercise}

\begin{exercise}[VI.29.E23: Noetherian warning]\label{exer:vi29-23}
Does Noetherian imply finite Krull dimension?
\end{exercise}

\begin{exercise}[VI.29.E24: Boundary VI/30]\label{exer:vi29-24}
Name one global dimension theorem deliberately deferred to VI/30.
\end{exercise}

\begin{exercise}[VI.29.E25: Boundary VI/31]\label{exer:vi29-25}
Which next chapter isolates codimension?
\end{exercise}

\begin{exercise}[VI.29.E26: Boundary VI/32]\label{exer:vi29-26}
Which next chapter treats local dimension geometrically?
\end{exercise}
